% Terminal note of the homily: everything the speech itself must not carry ---
% the audience, the spoken word count and the pace it implies, the relation to
% the three interpretations of the expansive study, the connections that are
% the homily's own, the translation each quoted text is given in, and the exact
% loci --- followed by the References. Nothing here is read aloud.

\section*{Note on Sources and Delivery}
\runninghead{Sources and Delivery}

{\footnotesize

\textbf{Audience and occasion.} An adult parish assembly, at the Mass
\latin{Dominica quinta quae superfuit post Epiphaniam}, II classis, of the
1962 \work{Missale Romanum}, editio typica, pp.~420--421, nos.~1749--1758,
said on Sunday 8~November 2026, the twenty-fourth Sunday after Pentecost. The
speech presumes that it follows the Gospel, and it speaks of the Offertory,
Secret, Communion and Postcommunion as still to come. It makes no use of the
Preface.

\textbf{Length and pace.} The spoken body runs to about 1,400 words (1,398 by
a plain count with comments and markup removed), which is about 10.8
to 11.7 minutes at an unhurried 120 to 130 words a minute, by arithmetic
alone; no one has delivered or timed it. It was read through in full,
silently, for sense and ease of speech.

\textbf{Relation to the reviewed interpretations.} The homily follows the
first of the expansive study's three interpretations, \textit{Both Grow until
the Harvest}, and every patristic reading it uses, as listed under the
sources below, is one that interpretation gives. The second and third
interpretations are left to the studies.
So is the disagreement over how far correction may go before the harvest:
Chrysostom's prohibition of killing (\work{Hom.\ in Mt.}\ 46.1), Augustine's
acceptance of exile and fines in his letter to Vincentius (\work{Ep.}~93.10),
and Aquinas's conclusions in the \work{Summa} (II-II q.~10 a.~8; q.~11
a.~3). From that letter the speech takes only Augustine's sentence that the
Church goes on correcting whom she can and bearing with whom she cannot until
the end. The Fathers name the cockle with different breadth: Jerome reads the
enemy's seed as the teachings of heretics, Chrysostom names heretics, and
Augustine bad Christians; the speech gives each his own word. They also name
the servants differently: Jerome takes them for angels (\work{In Matthaeum}
II, PL~26 col.~93, \latin{Servos patrisfamilias, ne alios accipias quam
angelos}), Aquinas for good men (\work{Super Mt.}\ c.~13, on 13:24--30),
Honorius of Autun for prelates (\work{Speculum Ecclesiae}, the sermon headed
\latin{Dominica tercia}, PL~172); the speech follows Augustine, who counts
himself and his hearers among them. Of the Lord's explanation (Mt
13:37--43), which the Mass does not read, the speech keeps the gloss that the
good seed are the children of the kingdom.

\textbf{Connections that are the homily's own.} These are drawn in the
preacher's voice and credited to no Father. The remark that no farmer expects
darnel to turn into wheat, and the conclusion that the householder waits for
the cockle's sake as well as the wheat's, are the preacher's, drawn from the
change Jerome and Chrysostom find in \latin{ne forte} and Augustine's address
to the tares. Paul's bearing with and forgiving one another (Col 3:13) is
taken as the servants' work while both grow (the first interpretation's own
inference), and Rabanus's patience in the householder is said to be the same
patience; neither Rabanus nor Paul mentions the other's text. That the prayer
for enemies takes in the cockle, and asks for the change Augustine asks of
the tares, applies Jerome's reading of Jer 29:7 to the parable; Jerome does
not mention the parable, and the speech does not say that the Introit's
promise of a hearing covers that prayer. That those who would clear the field
are sinners asking pardon is drawn from Augustine's answer at Ps 129:2 and
the assembly's own singing of the cry of Ps 129:1, with the Secret's ``our
sins''; Augustine's wish for the zealous servants is then applied to the
assembly. The anxiety for the crop that Chrysostom sees in the servants' haste
is said to be brought to God in the Collect; Aquinas's second affliction, the
sins of those who live among the saints, is said to be the wheat's beside the
cockle, and the conclusion that God, and not the servants, frees the wheat
from it is the preacher's, following the expansive study. Bede's condition
for being heard is made the speech's one response. The Postcommunion's pledge
and effect are read with the parable's barn.

\textbf{Sources, translations and loci.} Each source is listed once, with the
translation the speech uses; ``paraphrased'' means told in the speech's own
words. \emph{The Mass}, nos.~1749--1758: Introit, Jer 29:11--12, 14; Collect;
Epistle, Col 3:12--17; Gradual, Ps 43:8--9 (44:8--9); Alleluia, Ps 129:1--2
(130:1--2); Gospel, Mt 13:24--30; Offertory, Ps 129:1--2 (130:1--2); Secret;
Communion, Mk 11:24; Postcommunion. The chants and orations are quoted in the
English of the Lasance \work{New Roman Missal} (1945), a historical study aid
and not an approved liturgical translation of the 1962 Missal: the Introit
(\work{Supplementary Mass No.~2}, p.~744), the Gradual and Alleluia (p.~745),
and the Secret and Communion (p.~746). The Collect (under the Fifth Sunday
after the Epiphany, p.~196) and the Postcommunion are paraphrased.
\emph{Scripture} is quoted in the Douay--Rheims, Challoner revision, at Mt
13:27, 28, 29, 30, 38 and 43, Col 3:13, Jer 29:7 and Mk 11:25; the rest of Mt
13:24--28 and 37--39 is paraphrased, close to the Douay's, and so is
1~Cor 5:1; Paul's command at 1~Cor 5:13 is quoted in the expansive study's
English of the Latin in which Jerome cites it, under Jerome below. \emph{Augustine}, in the Nicene and
Post-Nicene English: \work{Sermo}~73 (NPNF Sermon~XXIII on the New Testament
lessons), n.~1, the Lord's words to the zealous servants with the wish ``That
time will come, may it only find you wheat!'', quoted; nn.~1--2, the tares as
bad Christians, paraphrased; n.~3, the address to the tares, paraphrased;
n.~4, ``ye are such as may be deceived,'' the high seats and the angels ``who
cannot make mistakes'', quoted, and ``We are but men, the reapers are the
angels'' and the good who groan to be few among many, paraphrased;
\work{Ep.}~93.34, ``correcting as many as she can, and bearing with those
whom she cannot correct'', quoted; \work{Enarr.\ in Ps.}\ 129.2, ``Who is it
then who crieth? \dots\ A sinner'', quoted. \emph{Jerome}, from the Latin:
\work{In Matthaeum} II, on 13:36--43, PL~26 cols.~93--94, where
\latin{Auferte malum de medio vestrum} (1~Cor 5:13) is quoted in the
expansive study's English and the rest paraphrased (\latin{lolium};
\latin{ne ubi quid ambiguum est, cito
sententiam proferamus; sed Deo judicii terminum reservemus}; \latin{non
suspicionem criminis, sed manifestum reatum}; \latin{datur locus
poenitentiae}); \work{In Hieremiam} V, on 29:5--9, PL~24 cols.~891--892,
quoted in the expansive study's English (\latin{orare praecepit pro inimicis
nostris \dots\ ut non simus nostra tantum salute contenti, sed inimicorum
quoque quaeramus salutem}). \emph{Chrysostom}, paraphrased from the Nicene and
Post-Nicene English: \work{Hom.\ in Mt.}\ 46.1 (the servants' anxiety for the
crop; many of the tares may change and become wheat); \work{Hom.\ in Col.}\ 8
(the complaint a petty matter; serious charges to be forgiven; to love).
\emph{Rabanus Maurus}, paraphrased from the Latin: \work{In Matthaeum} IV, on
13:24--30, PL~107 cols.~947--948 (\latin{non statim ad vindictam prorumpere,
sed magis ad virtutem patientiae recurrere debemus}; \latin{usque ad finem
vitae}). \emph{Bede}, quoted in the expansive study's English of the Migne
Latin: \work{In Marci evangelium expositio} III, on 11:22--25 (\latin{et
absque dubio impetrabunt quod petunt, si tamen in se peccantibus aliis primo
ipsi dimittunt}). \emph{Aquinas}, paraphrased from the Latin: \work{Super
Ps.}\ 43 n.~3 (\latin{spiritualiter, quando peccant illi qui habitant apud
eos, et ab his afflictionibus liberantur a Deo}).

\par}

\Needspace{14\baselineskip}
\section*{References}
\runninghead{References}

{\footnotesize
\begin{itemize}\setlength{\itemsep}{0.15em}\setlength{\parskip}{0pt}
\item \work{Missale Romanum ex decreto Sacrosancti Concilii Tridentini restitutum, Summorum Pontificum cura recognitum}, editio typica (Typis Polyglottis Vaticanis, 1962), CMAA facsimile, Proprium de Tempore, pp.~420--421, nos.~1749--1758.
\item Lasance, F.~X., and Walsh, F.~A., \work{The New Roman Missal in Latin and English}, revised edition (New York: Benziger Brothers, 1945), p.~196 (Collect) and pp.~744--746 (\work{Supplementary Mass No.~2}: Introit, Gradual, Alleluia, Secret, Communion).
\item \work{The Holy Bible Translated from the Latin Vulgate} (Douay--Rheims, Challoner revision), Project Gutenberg text: Jer 29:7; Mt 13:27--30, 38, 43; Mk 11:25; Col 3:13.
\item St Augustine, \work{Sermo} 73 (Sermon XXIII on the New Testament lessons), nn.~1--4, Nicene and Post-Nicene Fathers, 1st ser., vol.~6; \work{Epistula} 93 to Vincentius, nn.~10 and 34, Nicene and Post-Nicene Fathers, 1st ser., vol.~1; \work{Enarrationes in Psalmos} 129, n.~2, Nicene and Post-Nicene Fathers, 1st ser., vol.~8.
\item St Bede, \work{In Marci evangelium expositio} III, on Mk 11:22--25, Migne Latin text (Wikisource).
\item Honorius Augustodunensis (Honorius of Autun), \work{Speculum Ecclesiae}, the sermon on the cockle headed \latin{Dominica tercia}, \work{Patrologia Latina} 172.
\item St Jerome, \work{Commentarii in Matthaeum} II, on 13:36--43, \work{Patrologia Latina} 26 (Paris, 1845), cols.~93--94; \work{In Hieremiam} V, on 29:5--9, \work{Patrologia Latina} 24 (Paris, 1865), cols.~891--892.
\item St John Chrysostom, \work{Homilies on Matthew} 46, Nicene and Post-Nicene Fathers, 1st ser., vol.~10; \work{Homilies on Colossians} 8, Nicene and Post-Nicene Fathers, 1st ser., vol.~13.
\item St Rabanus Maurus, \work{Commentaria in Matthaeum} IV, on 13:24--30, \work{Patrologia Latina} 107 (Paris, 1864), cols.~947--948.
\item St Thomas Aquinas, \work{Super Psalmos}, on Ps 43, n.~3, Corpus Thomisticum text; \work{Super Evangelium S.~Matthaei lectura}, cap.~13, \work{Opera omnia}, vol.~3 (Venice: Bettinelli, 1745); \work{Summa theologiae} II-II q.~10 a.~8 and q.~11 a.~3, English Dominican translation; paraphrased or cited, not quoted.
\end{itemize}
\par}
